%%%%%%%%%%  scan idx 154  |  folio 141  %%%%%%%%%%
\SGAPrintPage{141}%
\sgaStableTarget{sga6.e01.para.6-2}6.2. Let \(C\) be a triangulated category and \(f\) an additive function from \(C\) to an abelian group \(F\). Then \(f\) induces an additive function on the additive category \(C^{\mathrm{ad}}\) underlying \(C\), hence an inclusion \(\Add(C,F)\subset\Add(C^{\mathrm{ad}},F)\). In particular, we have \(f(0)=0\) and \(f(L)=f(M)\) when \(L\simeq M\). On the other hand, we have \(f(L[n])=(-1)^n f(L)\) for every \(L\in\Ob(C)\) and \(n\in\mathbb Z\).

\sgaStableTarget{sga6.e01.para.6-3}6.3. Let \(C\) be an additive (resp.\ triangulated) category. The group \(\Add(C,F)\) depends functorially on \(F\). The functor \(\Add(C,.)\) is representable (cf.\ SGA 5 VIII 2) by an abelian group denoted by \(k(C)\)\SGAAuthorFootnote{raisedstar}{we shall not use here the usual notation \(\Kpx(C)\), since \(\Kpx(C)\) already denotes, when \(C\) is additive, the category of complexes of \(C\) ``up to homotopy''.} endowed with an additive map \(\cl_{C}:\Ob(C)\longrightarrow k(C)\) (when \(C\) is additive, we take for \(k(C)\) the symmetrization of the monoid of isomorphism classes of objects of \(C\); when \(C\) is triangulated, we take for \(k(C)\) the quotient of the free abelian group generated by \(\Ob(C)\) by the relations identifying \(L\) with \(L'+L''\) whenever there is a distinguished triangle \(L'\longrightarrow L\longrightarrow L''\longrightarrow L'[1]\)).

The group \(k(C)\) depends functorially on \(C\) (with respect to additive (resp.\ exact) functors). When \(u,u':C\to D\) are isomorphic additive (resp.\ exact) functors, we have \(k(u)=k(u')\).

When \(C\) is a triangulated category, we have a canonical epimorphism \(k(C^{\mathrm{ad}})\longrightarrow k(C)\) (cf.\ (6.2)).

\sgaStableTarget{sga6.e01.lem.6-4}\textbf{Lemma 6.4.} \uline{Let} \(C\) \uline{be an additive category}. \uline{The canonical functor} \(C\longrightarrow\Kpx^b(C)\) \uline{induces an isomorphism} (\uline{functorial in} \(C\)) \(k(C)\longrightarrow k(\Kpx^b(C))\).

\uline{Proof}. This is an immediate consequence of the following formula (which follows trivially from (6.2)):
\sgaStableTarget{sga6.e01.eq.6-4-1}\begin{equation}\tag{6.4.1}
  f(L)=\sum (-1)^i f(L^i)
\end{equation}
%%%%%%%%%%  scan idx 155  |  folio 142  %%%%%%%%%%
\SGAPrintPage{142}%
for \(L\in\Ob\Kpx^b(C)\) and \(f\) an additive function on \(\Kpx^b(C)\).

\sgaStableTarget{sga6.e01.def.6-5}\textbf{Definition 6.5.} \uline{Let} \(\underline{S}\) \uline{be a category,} \(\underline{C}\) \uline{an additive (resp}.\ \uline{triangulated)} \(\underline{S}\)\uline{-category} (1.1.1)\uline{, and} \(F\) \uline{an abelian presheaf on} \(\underline{S}\). \uline{A family of functions} \(f_X:\Ob(\underline{C}_X)\longrightarrow F(X)\) \uline{indexed by} \(X\in\Ob \underline{S}\) \uline{is called an additive} \(\underline{S}\)\uline{-function from} \(\underline{C}\) \uline{to} \(F\) \uline{when the following conditions are satisfied}:
\begin{enumerate}[label=(\roman*)]
\item \uline{for every} \(X\in\Ob \underline{S}\)\uline{,} \(f_X\) \uline{is an additive function} (6.1);
\item \uline{for every arrow} \(u:X\to Y\) \uline{of} \(\underline{S}\)\uline{, the following diagram is commutative}
\[
\begin{tikzcd}
\Ob\underline{C}_Y \arrow[r,"f_Y"] \arrow[d,"u^*"'] & F(Y) \arrow[d,"F(u)"]\\
\Ob\underline{C}_X \arrow[r,"f_X"] & F(X).
\end{tikzcd}
\]
\end{enumerate}
\uline{The additive} \(\underline{S}\)\uline{-functions from} \(\underline{C}\) \uline{to} \(F\) \uline{form an abelian group, denoted by} \(\Add_{\underline{S}}(\underline{C},F)\).

\sgaStableTarget{sga6.e01.rem.6-6}\textbf{Remarks 6.6.} a) Denote by \(\underline{C}^{\is}\) the fibered \(\underline{S}\)-category whose fiber over \(X\in\Ob \underline{S}\) is the category with the same objects as \(\underline{C}_X\), but whose arrows are the isomorphisms of \(\underline{C}_X\). On the other hand, denote by \(\underline F\) the fibered \(\underline{S}\)-category whose fiber over \(X\in\Ob \underline{S}\) is the discrete category (i.e.\ the only arrows are the identity arrows) \(F(X)\), the inverse-image functor by \(u\in\mathrm{Fl}(\underline{S})\) being \(F(u)\). Then an additive \(\underline{S}\)-function from \(\underline{C}\) to \(F\) is precisely a cartesian \(\underline{S}\)-functor from \(\underline{C}^{\is}\) to \(\underline F\), inducing an additive function in the sense of (6.1) on each fiber. It will be convenient for us to interpret additivity as follows. To fix ideas, suppose that \(\underline{C}\) is triangulated. The category \(\Cart_{\underline{S}}(\underline{C}^{\is},\underline F)\) (resp.\ \(\Cart_{\underline{S}}(\Tr(\underline{C})^{\is},\underline F)\) (1.1.2)) of cartesian \(\underline{S}\)-functors from \(\underline{C}^{\is}\) (resp.\ \(\Tr(\underline{C})^{\is}\)) to \(\underline F\) is the discrete category associated to an abelian
%%%%%%%%%%  scan idx 156  |  folio 143  %%%%%%%%%%
\SGAPrintPage{143}%
group, and we have a homomorphism
\[
  \phi:\Cart_{\underline{S}}(\underline{C}^{\is},\underline F)
  \longrightarrow
  \Cart_{\underline{S}}(\Tr(\underline{C})^{\is},\underline F),
\]
defined by \(\phi(f)(T)=f(L)-f(L')-f(L'')\) for \(f\in\Cart_{\underline{S}}(\underline{C}^{\is},\underline F)\) and \(T=(L'\to L\to L''\to L'[1])\in\Ob\Tr(\underline{C})\). To say that \(f\in\Cart_{\underline{S}}(\underline{C}^{\is},\underline F)\) is additive means that \(\phi(f)=0\); in other words, we have a (canonical functorial) exact sequence
\sgaStableTarget{sga6.e01.eq.6-6-1}\begin{equation}\tag{6.6.1}
  0\longrightarrow \Add_{\underline{S}}(\underline{C},F)
  \longrightarrow \Cart_{\underline{S}}(\underline{C}^{\is},\underline F)
  \xrightarrow{\phi}
  \Cart_{\underline{S}}(\Tr(\underline{C})^{\is},\underline F).
\end{equation}

b) Let \(u:X\to Y\) be an arrow of \(\underline{S}\). The ``inverse-image by \(u\)'' functors, being all isomorphic to one another, define (6.3) a unique homomorphism \(k(Y)\longrightarrow k(X)\). Consequently, \(X\mapsto k(X)\) defines an abelian presheaf on \(\underline{S}\), denoted by \(k(\underline{C})\). An additive \(\underline{S}\)-function from \(\underline{C}\) to \(F\) can then be interpreted simply as a homomorphism of presheaves \(k(\underline{C})\longrightarrow F\); in other words, we have
\[
  \Add_{\underline{S}}(\underline{C},F)=\Hom(k(\underline{C}),F).
\]
The presheaf \(k(\underline{C})\) of course depends functorially on \(\underline{C}\) with respect to additive (resp.\ exact) \(\underline{S}\)-functors.

\sgaStableTarget{sga6.e01.prop.6-7}\textbf{Proposition 6.7.} \uline{Let} \(\underline{S}\) \uline{be a site,} \(\underline{C}\) \uline{a flat abelian} \(\underline{S}\)\uline{-category, and} \(\underline{C}_0\) \uline{a strictly full additive sub-}\(\underline{S}\)\uline{-category of} \(\underline{C}\) \uline{satisfying the hypotheses} (4.0). \uline{Then the morphism of presheaves}
\[
  k(\underline{C}_0)\longrightarrow k(\Parf(\underline{C},\underline{C}_0))\qquad(\text{cf. }(4.9)),
\]
\uline{induced by the inclusion of} \(\underline{C}_0\) \uline{in} \(\underline{\Parf}(\underline{C},\underline{C}_0)\)\uline{, defines an isomorphism on the associated sheaves. In other words, for every abelian sheaf} \(F\) \uline{on} \(\underline{S}\)\uline{,}
%%%%%%%%%%  scan idx 157  |  folio 144  %%%%%%%%%%
\SGAPrintPage{144}%
\uline{the restriction homomorphism}
\[
  \Add_{\underline{S}}(\Parf(\underline{C},\underline{C}_0),F)
  \longrightarrow
  \Add_{\underline{S}}(\underline{C}_0,F)
\]
\uline{is an isomorphism}.

\uline{Proof}. Let \(F\) be an abelian sheaf on \(\underline{S}\); let us show that the arrow
\[
  \Add_{\underline{S}}(\Parf(\underline{C},\underline{C}_0),F)
  \longrightarrow
  \Add_{\underline{S}}(\underline{C}_0,F)
\]
is an isomorphism. It factors as
\[
\Add_{\underline{S}}(\Parf(\underline{C},\underline{C}_0),F)
 \xrightarrow{a}
\Add_{\underline{S}}(\Kpx^b(\underline{C}_0),F)
 \xrightarrow{b}
\Add_{\underline{S}}(\underline{C}_0,F).
\]
Now, by (6.4), the arrow \(b\) is an isomorphism. We are therefore reduced to proving that \(a\) is an isomorphism.

Consider the following commutative diagram, in which the rows are exact (6.6.1), and the vertical arrows are the ``restriction'' arrows:
\[
\begin{array}{ccccccccc}
0&\to&\Add_{\underline{S}}(\Parf(\underline{C}),F)&\to&\Cart_{\underline{S}}(\Parf(\underline{C})^{\is},\underline F)&\to&\Cart_{\underline{S}}(\Tr(\Parf(\underline{C}))^{\is},\underline F)\\
&&\downarrow a&&\downarrow&&\downarrow\\
0&\to&\Add_{\underline{S}}(\Kpx^b(\underline{C}_0),F)&\to&\Cart_{\underline{S}}(\Kpx^b(\underline{C}_0)^{\is},\underline F)&\to&\Cart_{\underline{S}}(\Tr(\Kpx^b(\underline{C}_0))^{\is},\underline F).
\end{array}
\]
Since \(F\) is a sheaf, \(\underline F\) is a stack (G II 2.2.4); thus, by (4.10), the two vertical arrows on the right are isomorphisms.

The same is therefore true of \(a\), q.e.d.

\sgaStableTarget{sga6.e01.rem.6-7-1}\textbf{Remark 6.7.1.} (cf.\ 4.10.1). In (6.7), we can replace \(\Parf(\underline{C},\underline{C}_0)\) by \(\Dpx^*(\underline{C})_{\underline{C}_0\text{-}\mathrm{parf}}\) (\(*=-\) or \(b\)).

%%%%%%%%%%  scan idx 158  |  folio 145  %%%%%%%%%%
\SGAPrintPage{145}%
\sgaStableTarget{sga6.e01.para.6-8}6.8. Let \((\underline{S},A)\) be a ringed topos and \(\underline{C}\) and \(\underline{C}_0\) the fibered \(\underline{S}\)-categories considered in (4.8). We have homomorphisms of abelian presheaves
\[
  k(\underline{C}_0)
  \xrightarrow{\alpha}
  k(\Dpx^-({}_A \underline{S})_{\mathrm{parf}})
  \xrightarrow{\beta}
  k(\Parf({}_A \underline{S})),
\]
which, by (6.7), induce isomorphisms on the associated sheaves. The homomorphism \(\alpha\) is compatible with inverse images, i.e.\ when \(u:(\underline{S}',A')\longrightarrow(\underline{S},A)\) is a morphism of ringed topoi, we have a commutative diagram, in which the vertical arrows are induced by the functor \(\mathbb L u^*\) (cf.\ (4.19.1)):
\[
\begin{tikzcd}
 k(\underline{C}_0) \arrow[r] \arrow[d] & k(\Dpx^-({}_A \underline{S})_{\mathrm{parf}}) \arrow[d]\\
 k(\underline{C}'_0) \arrow[r] & k(\Dpx^-({}_{A'} \underline{S}')_{\mathrm{parf}})
\end{tikzcd}
\]
(\(\underline{C}'_0\) being defined analogously to \(\underline{C}_0\)). When \(A\) is commutative, the tensor product \(\overset{\mathrm L}{\otimes}\) equips \(k(\underline{C}_0)\) and \(k(\Dpx^-({}_A \underline{S})_{\mathrm{parf}})\) with structures of presheaves of rings (cf.\ (4.19.1)), and \(\alpha\) is a homomorphism of presheaves of rings.

\sgaStableTarget{sga6.e01.para.6-9}6.9. In the situation of (6.8), suppose that \((\underline{S},A)\) is locally ringed (cf.\ (2.15.1)). Then (2.15.1), for every \(X\in\Ob \underline{S}\), \(\underline{C}_{0X}\) is the category of locally free \(A_X\)-Modules of finite type. Denote by \(\underline{C}_1\) the fibered \(\underline{S}\)-category defined by \(X\longmapsto\) (the category of free \(A_X\)-Modules of finite type). It is clear that \(\underline{C}_1\) satisfies the hypotheses (4.0). An object of \(\Dpx(\underline{S})\) is perfect if and only if it is perfect relative to \(\underline{C}_1\). We have homomorphisms of abelian presheaves
\[
  k(\underline{C}_1)\longrightarrow k(\underline{C}_0)\longrightarrow
  k(\Dpx^-(\underline{S})_{\mathrm{parf}})\longrightarrow k(\Parf(\underline{S}))
\]
